% \begin{figure}[t]
% \centering
% \includegraphics[width=\textwidth]{images/tweet_figure.png}
% \label{fig:languages_repos}
% \caption{\benchmarkname~is a dataset with challenging, enterprise-level, long-horizon software engineering tasks. Frontier models, such as GPT-5 and Claude Opus 4.1, score less than 25\% on \benchmarkname~with the SWE-Agent \cite{yang2024sweagent} scaffold. We design the dataset with contamination resistance, difficulty filtering, and human augmentation/verification.}
% \end{figure}

\section{Introduction}

Large Language Model (LLM) agents have been widely adopted in modern software development workflows. SWE-bench~\cite{jimenez2024swebench} and related works~\cite{yang2024swebenchmm, zan2024multiswe, yang2024sweagent, zhang2025swebench, swebenchverified} establish the task of issue resolution as a de-facto standard for assessing their capability and usefulness. In this setting, an agent is given an entire codebase, a task description (e.g., a bug report or feature request) in natural language and is instructed to produce a code patch that resolves the issue and passes the repository's test suite. These benchmarks have been instrumental in demonstrating both the substantial potential and the persistent limitations of current models as SWE agents.

% Notably, the state-of-the-art agents have reported over 70\% pass rate on SWE-Bench-Verified~\cite{swebenchverified}, a subset of SWE-Bench that is verifiably solvable by human programmers. In the next 6 - 12 months, there will be diminishing feedback from SWE-Bench-Verified to improve coding agents. Towards this end, this paper is motivated to (1) mitigate existing issues in SWE-Bench and (2) generate high-quality coding problems for evaluating the progress of LLM agents after SWE-Bench is saturated. As a result, we introduce  \benchmarkname. 

Current coding benchmarks face several limitations. First, many benchmarks are susceptible to \emph{contamination}~\cite{xu2024benchmark, deng2024investigating, Zhang2024ACE, White2024LiveBenchAC}, as exemplified by recent works~\cite{xu2024benchmark,deng2024investigating,cheng2025survey} and social media posts~\cite{aleithan2024swe,zhang2025swebench}. This risk arises because widely used open-source repositories—particularly those distributed under permissive licenses (e.g., MIT, Apache 2.0, BSD)—are prime candidates for inclusion in the large-scale web-crawled corpora used to pre-train LLMs~\cite{brown2020language}. As a result, constructing benchmarks from public GitHub repositories is inherently difficult, since many are already accessible as training data. Second, existing tasks may \emph{not} adequately capture the complexity of real-world software engineering. For example, SWE-Bench Verified~\cite{jimenez2024swebench} includes a substantial proportion of relatively trivial problems (161 out of 500) that require only one- to two-line modifications. In contrast, industrial software engineering, particularly in enterprise settings, often demands multi-file modifications spanning hundreds of lines~\cite{hassan2009predicting,steidl2017evaluating}. This discrepancy raises concerns about whether current benchmarks truly reflect the challenges faced in practical development scenarios.



% data contamination, lack of difficult, long-horizon tasks, and representation of real-world enterprise environments. 


% Recent studies have documented contamination concerns in popular benchmarks \cite{xu2024benchmark,deng2024investigating,cheng2025survey}, with practitioners reporting discrepancies between benchmark performance and real-world effectiveness \cite{aleithan2024swe,zhang2025swebench}.

% The first concern is the issue of \textbf{data contamination} \cite{xu2024benchmark,deng2024investigating, Zhang2024ACE, White2024LiveBenchAC}. The vast majority of popular open-source repositories, particularly those with permissive licenses (e.g., MIT, Apache 2.0, BSD), are prime candidates for inclusion in the massive web-crawled datasets used to pre-train LLMs \cite{brown2020language}. This makes it difficult to create a benchmark from public Github repos, as many of them are fair game for training data. 

Our first contribution in \benchmarkname~is a novel data collection strategy designed to \textbf{mitigate data contamination}. Specifically, our approach involves two complementary measures: (1) exclusively selecting repositories distributed under strong \emph{copyleft} licenses (GPL) to construct a public set (11 repositories) and a held-out set (12 repositories), and (2) \emph{acquiring commercial codebases} from real startups to capture enterprise-grade problems in a commercial set (18 repositories). In doing so, we reduce contamination risks by leveraging both legal protections and restricted data access. While analogous efforts may have been undertaken in industry using proprietary codebases, to the best of our knowledge, this work is the first to systematically apply such a methodology for curating a benchmark in the research community. The three subsets are made available under different access policies. The public set provides both problems and evaluation results openly. The held-out set remains private, preserving it for future overfitting checks against the public set. Finally, for the commercial set, we release evaluation results while keeping the underlying codebases private.

% Our first contribution in \benchmarkname~is its unique way of data collection to \textbf{mitigate data contamination}. Our approach includes (1) exclusively selecting repositories governed by strong \emph{copy-left} licenses (GPL) to form a public set (11 repos) and a held-out set (12 repos) and (2) \emph{acquiring commercial codebases} from real startups to source real-world enterprise problems in creating a commercial set (18 repos). Namely, we mitigate data contamination through legal barriers and private access restrictions. While similar efforts may have occurred in private industry with closed codebases, this work is, to our knowledge, the first to apply the approach to curate a benchmark in research publications. Problems and evaluation results from the public set are publicly available, while both are private for the held-out set, which will be later used to check for overfitting on the public set. For the commercial set, we will release the results and keep the codebases private.


% This results in 11 copy-left repositories in the public set, and 12 in a held-out set. Our dataset additionally includes 18 proprietary repositories in a commercial set, which are a first-of-it's-kind and purchased from private startups. 

 
% This creates a strong possibility that models evaluated on benchmarks derived from these repositories have been exposed to the solutions during training. Consequently, it becomes difficult to distinguish between genuine problem-solving ability and memorization, threatening evaluation validity. An ideal benchmark must provide verifiable guarantees that its contents are unlikely to be part of a model's training corpus.

% Another limitation is the homogeneity of problem domains and task complexity. Software engineering tasks are long-horizon tasks: a human might take days to weeks to complete the task, and an agent might take a few hours to complete the task. Existing benchmarks predominantly feature simple modifications: SWE-Bench tasks average 11.6 lines of code changes across 1.4 files \cite{jimenez2024swebench}, whereas industrial software engineering typically involves multi-file modifications spanning hundreds of lines \cite{hassan2009predicting,steidl2017evaluating}. 

The second contribution of \benchmarkname~is its emphasis on \textbf{challenging, diverse, and industrially relevant} tasks. To ensure task complexity, we exclude trivial edits (1–10 lines of code) and retain only problems requiring substantial, multi-file modifications. On average, the reference solutions span 107.4 lines of code across 4.1 files. Every problem involves at least 10 lines of change, and over 100 tasks demand more than 100 lines of modification. In addition to complexity, we prioritize diversity and representativeness. The curated repositories are all actively maintained and span a range of domains, including consumer applications, B2B services, and developer tooling platforms. Each repository contributes between 50 and 100 instances, with a strict cap of 100 instances, thereby reducing the risk of overfitting to any single repository. 

% The second contribution of \benchmarkname~is that we focus on \textbf{challenging, diverse and industrially-relevant} tasks. First, we filter to remove simple changes (between 1-10 lines of code) and focus on tasks that would require large, multi-file changes. That is, the reference solutions require 107.4 lines of code across 4.1 files on average. All problems require at least 10 lines of change, and more than 100 problems are in need of at least 100 lines of change. Second, repositories sourced here are all actively maintained across consumer applications, B2B services, and developer tooling platforms, with each repository contributing 50-100 instances to prevent overfitting. Since repositories are capped at 100 instances, models are tested on general capabilities rather than a focus on one particular repository.

The third contribution of \benchmarkname~is to demonstrate a \textbf{human-centered augmentation and verification workflow} to ensure task resolvability. We design a novel three-stage human-in-the-loop process that serves dual purposes: (1) clarifying ambiguity and adding missing context to preserve core technical challenges, and (2) recovering unit tests as robust verifiers by constraining solution spaces to avoid false negatives while maintaining implementation flexibility.

% In this work, we filter to remove simple changes (between 1-10 lines of code) and focus on evaluating agent performance on large, multi-file changes representative of the challenge of enterprise software engineering.

% Finally, current benchmark creation methodologies oversimplify problem-solving by filtering out issues deemed underspecified or ambiguous \cite{aleithan2024swe}. While this improves baseline quality, it creates an over-sanitized evaluation environment. In practice, software engineers frequently confront vaguely defined tasks requiring clarification and reasoned assumptions. Additionally, this filtering process inadvertently removes complex problems whose unit tests were written with specific implementations in mind, making them unsuitable for solution-agnostic evaluation without proper augmentation.

% To summarize, we establish several principles for \benchmarkname~to create a realistic and strong evaluation of software engineering agents in long-horizon, enterprise environments:

% \begin{itemize}
%     \item \textbf{Non-Contamination by Design}: We exclusively select repositories governed by strong copy-left licenses (GPL) and proprietary commercial codebases, providing robust safeguards against data contamination through legal barriers and private access restrictions. We include 23 such repositories -- 11 in the public set, and 12 in a held-out set. Our dataset additionally includes 18 proprietary repositories in a commercial set, which are a first-of-it's-kind and purchased from private startups.
    
%     \item \textbf{Diverse and Industrially-Relevant Tasks}: Our dataset spans 41 complex, actively maintained repositories across consumer applications, B2B services, and developer tooling platforms, with each repository contributing 50-100 instances to prevent overfitting. Since repositories are capped at 100 instances, models are tested on general capabilities rather than a focus on one particular repository.
    
%     \item \textbf{Balanced and Challenging Construction}: \benchmarkname~contains long-horizon tasks which would take hours to days for a software engineer to solve. We curate tasks requiring substantial modifications, with reference solutions averaging 107.4 lines of code across 4.1 files. We include 50+ problems that contain more than 10 lines of change, and 100+ problems with more than 100 lines of change, testing models on their ability to complete long-horizon tasks.
    
%     \item \textbf{Human-Augmented and Verified}: We employ a novel three-stage human-in-the-loop process that serves dual purposes: (1) clarifying ambiguity and adding missing context to preserve core technical challenges, and (2) recovering unit tests as robust verifiers by constraining solution spaces to avoid false negatives while maintaining implementation flexibility.
% \end{itemize}

% The human augmentation process deserves particular attention as it addresses a fundamental challenge in benchmark construction. Rather than discarding under-specified issues, our expert annotators systematically rewrite problem statements, add explicit requirements, and provide interface specifications. This methodology allows rigorous testing of reasoning and implementation capabilities while preserving the inherent difficulty of original tasks.

% \benchmarkname\ provides a contamination-resistant testbed with human-verified problems across three subsets: \benchmarkname~Public, \benchmarkname~Commercial, and an additional held-out private set. Our primary contributions are:

% \begin{itemize}
%     \item A large-scale benchmark for repository-level issue resolution with 1865 instances (731 public instances) across 41 repositories (11 in the public set, 12 in the held-out set, and 18 repos from enterprise startups). Notably, we are the first benchmark incorporating proprietary commercial codebases.
    
%     \item A framework for evaluating challenging software engineering tasks while ensuring resolvability: our tasks include the issue description as well as task requirements and interface modifications, resolving ambiguity issues for autonomous agent evaluation.
    
%     \item Comprehensive evaluation of state-of-the-art models establishing new baselines, where top models (GPT-5, Opus 4.1) achieve <25\% accuracy compared to >70\% on SWE-Bench Verified, revealing significant capability gaps in multi-file reasoning and complex codebase navigation.
% \end{itemize}

% We believe \benchmarkname\ will steer the community towards developing models that are not only powerful but also practical and aligned with real-world software development demands.

% Overall, LLM agents achieve only modest resolution rates on \benchmarkname~($\leq$23.3\% on the public set; $\leq$17.8\% on the commercial set), substantially lower than the $>$70\% Pass@1 reported on SWE-Bench Verified~\cite{swebenchverified}. We additionally observe a marked performance gap between the public and commercial sets, underscoring the greater complexity of enterprise codebases. Performance also varies systematically by programming language and repository: models generally perform better on Python and Go tasks, while several JavaScript/TypeScript repositories yield considerably lower results. To further characterize model behavior, we employ an LLM-as-a-judge analysis that surfaces distinct failure modes. Larger models (e.g., Opus 4.1) often fail on semantic or algorithmic correctness in large, multi-file edits, whereas smaller models (e.g., Qwen 3 32B) more frequently fail due to issues in syntax and formatting, tool use, or context management.

% Overall, LLM agents achieve modest resolve rates on \benchmarkname~($\leq$23.3\% on the public set; $\leq$17.8\% on the commercial set), far below the $>$70\% Pass@1 reported on SWE-Bench-Verified~\cite{swebenchverified}. We further observe a pronounced gap between the public and commercial sets, reflecting the added complexity of enterprise codebases. Performance varies systematically by language and repository (Python/Go generally higher; several JS/TS and specific repos lower). Finally, we employ an LLM-as-a-judge analysis that highlights distinct failure modes: some stronger models predominantly miss on semantic/algorithmic correctness for large, multi-file edits (e.g., Opus 4.1), whereas smaller models (e.g., Qwen 3 32B) more often fail due to syntax/formatting, tool-use, and context-management issues.


Taken together, \benchmarkname~aims to serve the community by providing a contamination-resistant and industrially realistic benchmark, supported by a transparent curation process and fine-grained diagnostic analyses. We release both the problems and evaluation results for the public set, retain the held-out set to monitor potential overfitting, and report results on the commercial set while preserving the privacy of its underlying codebases. Combined with standardized evaluation protocols and trajectory-level failure analyses, \benchmarkname~offers a rigorous foundation for measuring progress beyond the saturation of SWE-Bench Verified, establishing a common yardstick for researchers and practitioners developing next-generation coding agents.

% Taken together, \benchmarkname~aims to serve the community by providing a contamination-resistant, industrially realistic benchmark with a transparent curation process and fine-grained diagnostics. We release the public set and its evaluation results, retain the held-out set for overfitting checks, and report results on the commercial set while keeping codebases private. Together with standardized evaluation settings and trajectory-level failure analyses, our work offers a high-signal basis for measuring progress beyond SWE-Bench-Verified saturation and a common yardstick for both researchers and practitioners building next-generation coding agents.


